\documentclass[10pt,a4paper]{article}
\usepackage[utf8]{inputenc}
\usepackage[T1]{fontenc}
\usepackage[spanish]{babel}
\usepackage[margin=2.5cm]{geometry}
\usepackage{amsmath,amssymb,booktabs,tabularx}
\usepackage[hidelinks]{hyperref}
\hypersetup{pdftitle={Busqueda de metodos MILP con presupuestos de precision}}
\setlength{\parindent}{0pt}
\setlength{\parskip}{5pt}
\title{Búsqueda de un método para MILP\\con presupuestos de precisión}
\author{Nota de investigación separada del manuscrito OMS}
\date{17 de octubre de 2026}
\begin{document}
\maketitle

\textbf{Resultado.} Se desarrollaron y contrastaron cuatro candidatos, sin
restringir la búsqueda al escalado conjunto de filas y columnas. Dos tienen
propiedades matemáticas implementadas y verificadas; ninguno acredita todavía
una ventaja práctica propia frente a controles fuertes. No se añadió contenido
al artículo ni se modificó su versión congelada v1.3.14.

\section*{Resultados de la búsqueda}
\begin{tabularx}{\textwidth}{@{}p{3.8cm}X@{}}
\toprule
Candidato & Evidencia y decisión\\\midrule
Cambio mínimo de columnas & 27 casos; cambio cero en los 22 compatibles;
cinco certificados de incompatibilidad. No aparece oportunidad real.\\
Contracción certificada de big-$M$ & 1.753 deducciones de cotas y 96
contracciones de $10^{17}$ a aproximadamente $2{,}27\times10^6$.
Los contratos redondeados siguen incompatibles.\\
Escala de actividades SOS1 & 54 llamadas. La propuesta y SOS convencional
cumplen los presupuestos en 12/12 cada uno. El control obtiene mejores
objetivos en control30.\\
Copias enteras redundantes & 78 llamadas. Tres celdas admisibles y una
abstención; presolve elimina todas las copias. Los controles continuos
producen el mismo modelo reducido.\\\bottomrule
\end{tabularx}

Se conservaron y verificaron los 132 puntos devueltos. Se examinaron 3.456
conjuntos de hipótesis SOS y 2.592 de amplificación. Se hicieron además seis
diagnósticos de presolve sin optimizar, repetidos una vez para comprobar su
reproducción (12 llamadas en total). No se infiere optimalidad MILP de los
estados del solver ni de las comprobaciones primales.

\section*{La propiedad candidata de las copias enteras}
Para una binaria $z$, se añade una entera $q\in[0,K]$ y
$\sigma(Kz-q)=0$, con $K$ entero positivo. Esta extensión preserva la
proyección LP y MILP del modelo de partida. Si
\[
 |z-\operatorname{round}(z)|\leq\gamma,\qquad
 |q-\operatorname{round}(q)|\leq\gamma,\qquad
 |\sigma(Kz-q)|\leq\epsilon,
\]
y $K\gamma+\epsilon/\sigma+\gamma<1/2$, las dos redondeadas enteras cumplen
$\operatorname{round}(q)=K\operatorname{round}(z)$. En consecuencia,
\[
 |z-\operatorname{round}(z)|\leq
 \gamma_{\rm eff}=\frac{\gamma+\epsilon/\sigma}{K}.
\]
El límite para una fila con una sola binaria pasa a ser
$\epsilon/\alpha+|a_z|\gamma_{\rm eff}\leq\delta$.
El selector elige factores diádicos, comprueba exportación exacta y minimiza
el tamaño máximo de los coeficientes dentro de la ventana declarada.
Las hipótesis se verifican sobre los puntos almacenados.

En cinca, se admite $K=2048$ para presupuesto $10^{-4}$ y $K=2097152$
para $10^{-8}$. En control30, $10^{-4}$ admite $K$ entre 2 y 256; para
$10^{-8}$ el selector se abstiene. La integralidad redundante es una operación
conocida y no se afirma prioridad del criterio propuesto.

\textbf{Control causal.} Con presolve automático, las nueve parejas de
copias enteras y continuas producen exactamente el mismo punto completo y
objetivo. Los tres pares de matrices y dominios reducidos también son
idénticos. Se eliminan las 24 copias de control30 y las 96 de cada variante
de cinca. El mayor cociente de trabajo entero/continuo es aproximadamente
1,000001. Con presolve desactivado, en control30 la mediana de objetivo es
3,16534 con copias enteras y 2,66253 con continuas; menor es mejor.
No aparece una ventaja consistente del mecanismo.

\section*{El control SOS convencional también importa}
Para $g(x_C)\leq Mz$, la reformulación con actividad $y$ usa
$g-hy\leq0$, $0\leq y\leq M'/h$, $w+z=1$ y SOS1$(w,y)$, donde $M'$
es una contracción certificada común a los controles. El límite condicional
del error redondeado es
\[
 \frac{\epsilon}{\alpha}+h\max(\gamma_{\rm SOS},\tau),
\]
con $\tau$ el error permitido de la cota superior de $y$.
Escoger $h$ por ese contrato no mejoró SOS con $h=1$ en el piloto:

\begin{center}
\begin{tabular}{lrr}
\toprule
Presupuesto en control30 & Propuesta & SOS con $h=1$\\\midrule
$10^{-4}$ & 2,36943 & 1,83523\\
$10^{-8}$ & 5,45573 & 1,83523\\\bottomrule
\end{tabular}
\end{center}
Son medianas de objetivos de minimización al mismo límite de trabajo.
En cinca no aparece una diferencia relevante. Veinte conjuntos de hipótesis
nominales fallan en la propuesta de control30, aunque sus puntos redondeados
cumplen por comprobación directa. Se mantienen separadas esas dos evidencias.

\section*{Alcance y reproducción}
Los modelos son control30-3-2-3 y neos-3421095-cinca, seleccionados antes por
su diagnóstico. Son datos de descubrimiento. Los presupuestos $10^{-4}$ y
$10^{-8}$ son algebraicos experimentales y solo cambian en las activaciones.
Las otras filas y cotas conservan sus presupuestos originales.

Gurobi 13.0.2; cuatro hilos; WorkLimit=1, TimeLimit=30 s, MIPGap=0.01;
$\epsilon=10^{-6}$, $\gamma=10^{-9}$. Semillas SOS: 151, 157, 163;
semillas de copias: 167, 173, 179. Muchos puntos útiles terminan por límite
de trabajo; sus objetivos no son óptimos certificados. El tiempo de solver
no comprende toda la selección y exportación. OutputFlag=0 dejó únicamente
cabeceras en los logs; se preservan modelos, parámetros, atributos y puntos,
pero no trazas completas del árbol.

El informe íntegro está en \texttt{RESULTADOS.md}; las 58 agregaciones de
tres semillas, en \texttt{SUMMARY.csv}. Desde la carpeta de esta nota,
\texttt{python3 verify\_all.py} repite ocho verificaciones sin optimizar;
\texttt{python3 analyze.py} reconstruye las agregaciones. Se requiere el
repositorio completo. Cada candidato conserva protocolo, implementación,
pruebas y resultados. Los pilotos están protegidos contra sobrescritura.

\section*{Antecedentes que delimitan la novedad}
SCIP ya implementa indicadores y escalado de holguras en
\href{https://raw.githubusercontent.com/scipopt/scip/master/src/scip/cons_indicator.c}{\texttt{cons\_indicator.c}}.
La semántica numérica de SOS y sus conversiones está documentada por
\href{https://docs.gurobi.com/projects/optimizer/en/12.0/concepts/modeling/constraints.html}{Gurobi}.
La agregación de variables enteras forma parte del
\href{https://github.com/ERGO-Code/HiGHS/blob/master/highs/presolve/HPresolve.cpp}{presolve de HiGHS}.
El fortalecimiento de coeficientes es un antecedente establecido, entre otros,
en \href{https://research.dial.uclouvain.be/entities/publication/03ea6bfe-7105-4d5c-a48a-47578c173bd1}{Andersen y Pochet}.
Esta consulta inicial no prueba prioridad ni reemplaza una revisión exhaustiva.

\textbf{Decisión para OMS.} Mantener estos resultados como investigación
separada. Hay propiedades verificables y prototipos reproducibles, pero no
evidencia suficiente para anunciar un método nuevo con utilidad adicional.
\end{document}
